% ACHTUNG !!!!!!
% In diesem Dokument müssen Sie keine Anpassungen vornehmen. Nutzen Sie die DECKBLATT.tex um das Deckblatt auszufüllen!

\documentclass[a4paper, 12pt]{article}
\usepackage[ngerman]{babel}
\usepackage[utf8]{inputenc}
\usepackage[T1]{fontenc}
\usepackage{amssymb}
\usepackage{xcolor}
\usepackage{tikz}
\usepackage{tabularx}
\usepackage{enumitem}
\usepackage{listings}
\usepackage{graphicx}
\usepackage{hyperref}
\usepackage[top=2.5cm, bottom=2cm, left=2cm, right=2cm]{geometry}
\usepackage{fancyhdr}

\newcommand*{\defProjectLead}[2]{\def \projectLeadFirstname {#1} \def \projectLeadSurname {#2}}
\newcommand*{\defDeveloperZero}[2]{\def \developerZeroFirstname {#1} \def \developerZeroSurname {#2}}
\newcommand*{\defDeveloperOne}[2]{\def \developerOneFirstname {#1} \def \developerOneSurname {#2}}
\newcommand*{\defDeveloperTwo}[2]{\def \developerTwoFirstname {#1} \def \developerTwoSurname {#2}}
\newcommand*{\defDeveloperThree}[2]{\def \developerThreeFirstname {#1} \def \developerThreeSurname {#2}}

% Geben Sie Ihrem Projekt einen Namen und füllen Sie aus, an welchem Datum Sie die Spezifikation einreichen.
\def \projectName {Similar Slides}
\def \submissionDate{28.11.2026}

% Hier muss der Teamleiter eingetragen werden. NUR Teamleiter laden am Ende die fertige PDF zur Abgabe hoch.
\defProjectLead{Niklas}{Kaesehagen}

% Hier bitte die übrigen Entiwckler/innen eintragen. Wenn weniger als 5, die übrigen einfach leer lassen.
\defDeveloperZero{Quentin}{Becker}
\defDeveloperOne{VORNAME}{NACHNAME}
\defDeveloperTwo{}{}
\defDeveloperThree{}{}

% Grundsätzliche sollte es einen Customer Proxy geben und der Rest sind Entwickler/innen. Wenn Sie OneTutor als Customer Proxy verwenden, tragen Sie alle Mitglieder als Entwickler/in ein.
\def \roleDeveloperZero{Entwickler*in}
\def \roleDeveloperOne{Entwickler*in}
\def \roleDeveloperTwo{}
\def \roleDeveloperThree{}

% ---- HINWEIS ----
% Datum und Unterschrift des Kunden sowie die Bewertung erfolgen durch Ihren Tutor.

\newcommand*{\swtdate}[0]{\the\year/\the\numexpr\year+1}
\newcommand*{\swtdateshort}[0]{\the\year}

\setlength{\headheight}{14.5pt}

% header style
\pagestyle{fancy}
\fancyhf{}
\rhead{\thepage}
\lhead{\textit{SWT \swtdate}}
\chead{\textit{\projectName}}

% pdf options
\hypersetup{
	pdfauthor = {Gruppe \projectName},
	pdftitle = {Spezifikation \projectName},
	pdfsubject = {Softwareprojekt},
	pdfcreator = {VORNAME NACHNAME}}

% command to add the signateure-fields
\newcommand*{\SignatureAndDate}[2]{%
    #2
    \par\noindent\makebox[\linewidth]{\hrulefill}%
    \par\noindent\makebox[\linewidth][l]{#1}%
}%

% command to add a TODO Line
\newcommand*{\todo}[1]{\textbf{\textcolor{red}{TODO: #1}\\}}

\definecolor{lightlightgray}{RGB}{222,222,222}

\begin{document}
	\thispagestyle{empty} % no header on first page
	\begin{center}
		{\Huge Softwaretechnik Wintersemester \swtdate}\\[.35cm]
		{\large am Fachgebiet Software Engineering, Leibniz Universität Hannover}\\[.1cm]
	\end{center}
	\begin{center}
	\colorbox{lightlightgray}{
		\begin{minipage}{\linewidth}
			\begin{center}
				\vspace{0.2cm}
				{\Huge Spezifikation}\\[.1cm]
				{\Huge \projectName}\\[.3cm]
				\texttt{SWT-WS\swtdateshort{}-{\projectLeadSurname}-Spec.pdf}
				\vspace{0.2cm}
			\end{center}
		\end{minipage}
	}
	\end{center}
	\textbf{Vorgelegt}\\
		\hspace*{1cm}am \submissionDate \\
		\hspace*{1cm}von \projectName\\[0.2cm]
	\textbf{Ausführende:}
		\begin{center}
			\begin{tabular}{|l|l|l|}
				\hline
				\textit{Nachname} 		& \textit{Vorname}          & \textit{Rolle} \\ \hline\hline
				\projectLeadSurname		& \projectLeadFirstname		& Teamleiter/in\\ \hline
				\developerZeroSurname	& \developerZeroFirstname	& \roleDeveloperZero\\ \hline
				\developerOneSurname	& \developerOneFirstname	& \roleDeveloperOne\\ \hline
				\developerTwoSurname	& \developerTwoFirstname	& \roleDeveloperTwo\\ \hline
				\developerThreeSurname	& \developerThreeFirstname	& \roleDeveloperThree\\ \hline
			\end{tabular}
		\end{center}
		\vspace{0.3cm}
	\textbf{Das Dokument enthält}\\[0.1cm]
		\hspace*{1cm}$\boxtimes$ Die Anforderungen aus Kundensicht (User Requirements)\\[0.1cm]
		\hspace*{1cm}$\boxtimes$ Anforderungen, wie das zu System zu gestalten ist (System Requirements)\\[1cm]
	\SignatureAndDate{Datum, Unterschrift des Projektleiters, auch für die anderen Projektangehörigen}{\submissionDate}\\[0.3cm]
	\textbf{Kunden-Bewertung}\\[0.4cm]
		Der Kunde \underline{~~~~~~~~~~~~~~~~~~~~}, bestätigt mit seiner Unterschrift, diese Anforderungsspezifikation erhalten, geprüft und für inhaltlich $\square$ \textbf{in Ordnung} \textbar~ $\square$ \textbf{weitgehend in Ordnung} \textbar ~ $\square$ \textbf{deutlich zu verbessernd} \textbar ~ $\square$ \textbf{nicht akzeptabel} befunden zu haben.\\[1cm]
	\SignatureAndDate{Datum, Unterschrift des Kunden; evtl. Vermerk.}{}

	\pagebreak
	\tableofcontents
	\pagebreak

	\raggedright

	\section{Mission des Projekts}

\subsection{Erläuterung des zu lösenden Problems}
In jedem Veranstaltungszyklus entstehen hunderte neue Folien,  die ständig angepasst, optimiert und verändet werden. Dort den Überblick zu behalten und effektiv zu arbeiten gestaltet sich häufig als schwierig und erschwert den stressigen Vorlesungsbetrieb weiter.
Die Lösung dafür ist \textbf{Similar Slides}, ein zentrales Repository in dem alle Versionen einer bestimmten .pptx-Datei gesammelt, sortiert und versioniert werden. So kann sichergestellt werden, dass alle Mitarbeitenden einer Veranstaltung immer Zugriff auf die neuste und vergangene Versionen der Folien haben und sich die Unteschiede dieser schnell und klar erkennen lassen.
% Beschreiben Sie in diesem Paragraphen, welche Probleme durch das Projekt gelöst werden sollen. Dies dient als eine erste Einleitung und soll den Lesern ein Grundverständnis von Zweck und Einsatz der Software vermitteln. Einleitungen wie diese schreiben sich am besten, wenn bereits alle anderen Kapitel fertiggestellt wurden, da es sich hierbei um eine leicht vertändliche und kurze Zusammenfassung der folgenden Kapitel handeln soll. Orientieren Sie sich für diesen Abschnitt grob an folgender Struktur:

%   1. Wie heißt das Projekt
%   2. Beschreiben Sie das Problem (Welches Problem? Wann tritt es auf? In welchem Kontext?)
%   3. Beschreiben Sie, wie das Projekt das Problem lösen soll
%   4. Nennen Sie die möglichen Zielgruppen und wie diese von der fertigen Software beeinflusst werden

% Ca. 9-12 Zeilen Text wären hier angebracht.

\subsection{Wünsche und Prioritäten des Kunden}
	Es folgen die Wünsche des Kunden, nach absteigender Priorität geordnet:
	\begin{enumerate}
		\item Es muss eine Repository geben, in das die Foliensätze hochgeladen werden können
        \item Es sollen keine Folien verloren gehen
        \item Foliensätze, die genau so wie gespeichert erneut hochgeladen werden, sollen nicht gespeichert werden
        \item Leicht veränderte Folien sollen nebeneinander gestellt und nach Änderungsdatum sortiert gespeichert werden
        \item Einzelne Folien sollen aus dem Repo in PowerPoint hineinkopiert werden können
        \item Im Repo sollen kleiner Behälter angelegt werden können, in denen Mitarbeitende ihre Folien hochladen können
        \item Die Ähnlichkeit der Folien soll mit einem Ähnlichkeitsmaß gemessen werden
	\end{enumerate}

% nur funktionale anforderungen oder alle?

% Listen Sie hier die vom Kunden genannten Wünsche zum Endprodukt auf und geben Sie an, wie wichtig dem Kunden die verschiedenen Wünsche sind. Eine Sortierung würde sich hier anbieten.

\subsection{Domänenbeschreibung}

Similar Slides ist hauptsächlich auf den Lehrbetrieb einer Hochschule oder Universität ausgelegt. Dort soll es im Veranstaltungs- beziehungsweise Vorlesungsbetrieb unter Microsoft Windows in Kombination mit der PowerPoint Software eingesetzt werden. Dies geschieht während der Ausarbeitung der Folien, also in der Vorbereitung auf die tatsächliche Präsentation. 

% Erläutern Sie, in welcher Domäne das Endprodukt eingesetzt wird. Dazu gehören der Anwendungskontext (wo wird die Software eingesetzt? Z.B. auf dem Hochschulinformationstag) und der technische Kontext (Auf welchen Geräten läuft die Software? Z.B. Mobil, Tablet, Server). Gehen Sie entsprechend darauf ein, welche Effekte die Umgebung auf die Anwendung der Software haben könnte (Lichtverhältnisse, Textgröße, etc.)

% Ca. 8 Zeilen Text wären hier angebracht.
	
\subsection{Maßnahmen zur Anforderungsanalyse}
Zur Erhebungs der Spezifikationen wurde das Video des Kundeninterviews auf dessen Anforderungen untersucht. Anschließend wurden diese in Decision Records festgehalten. Alle offenen Fragen wurden mithilfe des Customer-Proxys geklärt und die Decision Records entsprechend aktualisiert.

% Bitte einfügen: Welche Mockups, Überlegungen, Prototypen wurden erstellt, um die Anforderungen zu erheben? Jeweils nur sehr kurz erläutern, wozu sie dienten, was dabei herausgekommen ist. Sie müssen an dieser Stelle KEINE Mockups etc. erstellen. Überlegen Sie sich, welche Maßnahmen HÄTTEN ergriffen werden können/müssen, um die Anforderungen erfolgreich zu erheben und welchen Zweck diese erfüllt hätten. Hier können auch weniger formale Ansätze gewählt werden, es müssen nicht zwingend die o.g. Mockups, Prototypen sein.
	\pagebreak
	\section{Rahmenbedingungen und Umfeld}

\subsection{Einschränkungen und Vorgaben}
Similar Slides soll in einer objektorientierten Programmiersprache entwickelt werden. Ferner soll \textit{ausschließlich} das proprietäre .pptx-Format unterstützen. Die Analyse und der Vergleich des Inhalts der Folien soll über einen Zugriff auf die XML-Struktur einer .pptx-Datei mithilfe von VisualBasic erfolgen. Das Ähnlichkeitsmaß vergleicht die Elemente einer Folie nach der folgenden Priorität: Zuerst der Titel, dann Graphiken und Bilder und zuletzt wird der tatsächliche Textinhalt beachtet.

% Erklären Sie hier, welche Vorgaben der Kunde bezüglich der Umsetzung des Projektes ausgesprochen hat. Wird beispielsweise verlangt, dass ein bestimmtest Framework oder eine Arbeitsweise vorausgesetzt werden? Müssen rechtliche Vorgaben eingehalten werden? Möglicherweise gibt es implizite Einschränkungen, wie eine Internetanbindung, Erreichbarkeit der Software oder die Verwendung von bestimmten Zahlungsdienstleistern.

% Mind. 6 Zeilen wären hier angebracht.

\subsection{Anwender}
Die Software richtet sich hauptsächlich an den Dozenten einer Veranstaltung und seine wissenschaftlichen Mitarbeiter. Aus dieser Personengruppe sollen besonders diese Angesprochen werden, die eine solide Basis an Office/PowerPoint-Erfahung vorzuweisen haben, aber wenig Zeit bei der Vorlesungs- und Übungsvorbereitung haben. 

% Beschreiben Sie hier die Mitglieder Ihrer Zielgruppe. Richtet sich die Software an Schüler, Studenten oder vielleicht Professoren? Beschreiben Sie für jedes Mitglied der Zielgruppe welche Charakteristiken diese ausmachen und wie dies die Nutzung des Endprodukts beeinflusst. Richtet sich beispielsweise Ihre Software an ältere Menschen könnte dies die Computerkompetenzen oder die Erkennung von kleinen Details auf einem Bildschirm beeinflussen.

% Mind. 6 Zeilen wären hier angebracht.

\subsection{Schnittstellen und angrenzende Systeme} 
Similar Slides muss mit dem PowerPoint-Programm interagieren, um das gezielte Einfügen einzelner Folien zu ermöglichen. Außerdem muss die Software den XML-Unterbau einer .pptx-Datei ausgelesen und analysiert können. Zuletzt muss auch noch auf ein zentrales Repository mit Lese- und Schreibvorgängen zugregriffen werden. 

% Nennen und erläutern Sie hier die verschiedenen Systeme / Libraries, die mit Ihrem Projekt interagieren müssen. Beipsielsweise verschiedene Drittanbieter, die Services wie Zahlungsabwicklung bereitstellen oder APIs, über die Sie Informationen beziehen oder Befehle ausführen müssen.
	\pagebreak
	\section{Funktionale Anforderungen}

% Fertigen Sie hier eine Liste aller funktionalen Anforderungen an. Halten Sie sich dabei an die Formatierung aus der Vorlesung und Übung. Sortieren Sie die Anforderungen nach Priorität. Die Priorität erhalten Sie aus dem Video. Falls etwas unklar ist, fragen Sie Ihren Customer-Proxy.

\subsection{Use Case-Diagramm}

% Viele der folgenden Unterkapitel sind leer. Füllen Sie sie entsprechend Ihren SWT-Kenntnissen; bedenken Sie dabei auch, dass ein UML-Diagramm oder auch ein Use Case ohne Erläuterung oft nicht aussagekräftig ist. Scheuen Sie sich also nicht, noch Kommentare oder weitere Angaben zu machen, wenn das der Klärung dient

\pagebreak

\subsection{Use Case-Beschreibungen}

% Abläufe mit interessanter Interaktion zwischen einem Akteur und verschiedenen Systemteilen werden als Use Cases beschrieben. Zu jedem Use Case kann es weitere, detaillierende Anforderungen geben, die dort angegeben sind (z.B R412). Rein interne Abläufe werden unter Umständen anders notiert (z.B. Ablauf- oder Sequenzdiagramm). Hier handelt es sich um Anforderungen des Kunden.


%%%% BEGIN - VORLAGE FÜR USE CASES %%%%
% Erstellen Sie eine eigene Datei für jeden Use Case

\subsubsection{UC: TITEL}

\begin{tabularx}{\linewidth}{|p{3.5cm}|X|}
	\hline
	\textbf{UC <01>}		& \textbf{<Titel>} \\ \hline
	\textbf{Umfeld}			& ... \\ \hline
	\textbf{Systemgrenze}	& ... \\ \hline
	\textbf{Ebene}			& ... \\ \hline
	\textbf{Hauptakteur}	& ... \\ \hline
	\textbf{Stakeholder}
        & Stakeholder1 möchte ...  \\
        & Stakeholder2 möchte ... \\
        & Stakeholder3 möchte ... \\
        \hline
	\textbf{Voraussetzungen}   & ... \\ \hline
	\textbf{Garantie}          & ... \\ \hline
	\textbf{Erfolgsfall}       & ... \\ \hline
	\textbf{Auslöser}          & ... \\ \hline	
	\textbf{Hauptszenario}			
        & 1. \\
	    & 2. \\
		& 3. \\
		& ... \\
        \hline
	\textbf{Erweiterung}			
        & 1a. WENN ... DANN ... \\ 
		& 2a. WENN ... DANN ... \\
        \hline
	\textbf{Priorität}     & ... \\ \hline
	\textbf{Technologie}   & ... \\ \hline
\end{tabularx}

\\[0.5cm]
\textbf{Erläuterungen und Details}
\begin{itemize}
    \item 
\end{itemize}

%%%% END - VORLAGE FÜR USE CASES %%%%
	\pagebreak
	\section{Qualitätsanforderungen}

% Fertigen Sie hier eine Liste aller nicht-funktionalen Anforderungen an.
\begin{enumerate}
    \item Der Code soll lesbar und gut modifizierbar sein
    \item Die Ähnlichkeit wird in der Priorität: Titel, Graphiken, Bilder, dann Text bestimmt
    \item Es soll über VisualBasic auf die XML-Struktur der Präsentationen zugegriffen werden
\end{enumerate}

\subsection{Qualitätsziele des Projekts}

% Beschreiben Sie hier die einzelnen Qualitätsziele, die an die Software gestellt werden und wie diese am Ende des Projektes gemessen werden können. Beispielsweise könnte Sicherheit ein Qualitätsziel sein. Konkretiesieren Sie dieses und beschreiben Sie, was (hier Sicherheit) in Ihrer Software bedeutet und an welchen Stellen Sie besonders Acht geben müssen. Um zu beweisen, dass die Qualitätsziele erreicht wurden, müssen diese gemessen werden können. Überlegen Sie sich ein Verfahren, um die einzelnen Ziele messen zu können (bei Sicherheit könnte dies Beispielsweise ein Pentest sein.)

% Hier sind nur grobe Überlegungen notwendig. Tierfergehend werden Sie dies erst in SWQ kennenlernen.

\subsection{Qualitäts-Prioritäten des Kunden}

% Anhand der Informationen aus dem Video sollen die Anforderungen hier priorisiert werden. Wenn etwas unklar ist, fragen Sie Ihren Customer-Proxy

Die Qualitätsziele sind wie folgt absteigend priorisiert:
\begin{itemize}
    \item ...
\end{itemize}

\subsection{Wie Qualitätsziele erreicht werden sollen}

	\pagebreak
	\section{Hinweise zur Umsetzung}

% Hier sollen alle weiteren, wichtigen Angaben stehen, wie zum Beispiel: Oberflächen-Mockups. Bitte unbedingt mit kurzer Erläuterung, was man da sieht und was Sie mit dem Mockup ausdrücken wollen!

% Sie können hier auch noch ganz andere Sachen einbauen, die Sie für wichtig halten. Das Dokument soll aber nicht unnötig aufgebläht werden. Kopieren Sie nicht irgendwelche allgemeinen Texte hinein.
	\pagebreak
	\section{Probleme und Risiken}
\begin{enumerate}
	\item \textbf{WENN} Die Unterscheidung von Folien aufgrund ihrer XML-Struktur nicht wie geplant funktioniert\\
	      \textbf{DANN} Ist die zentrale Anforderung nach einer einfachen und strukturierten Differenzierung der Folien nicht mehr erfüllt. Konsequenz: Es brauch eine neue schnelle und verlässliche alternative Methode zur Erkennung von Unterschieden zwischen Folien.

	\item \textbf{WENN} ...\\
	      \textbf{DANN} ... Konsequenz: ...
\end{enumerate}

\textbf{Abhilfe:}
\begin{enumerate}
	\item ...
	\item ...
\end{enumerate}

	\pagebreak
	\section{Optionen zur Aufwandsreduktion}

% Sehr wichtig: Klären Sie gleich zu Anfang, welche Funktionen und Anforderungen für den Kunden möglicherweise nicht ganz so wichtig sind. Dort kann man angreifen, wenn nicht alles zu machen ist. Falls Abstriche nicht klar sind, fragen Sie Ihren Customer-Proxy.

\subsection{Mögliche Abstriche}

\subsection{Inkrementelle Arbeit}
    


    
	\pagebreak
	\section{Glossar}

% Beschreiben Sie hier alle möglichen Fach- oder Domänenbegriffe, damit alle Stakeholder die Spezifikation nachvollziehen können.

\begin{description}
    \item[KEYWORD] DESCRIPTION
\end{description}

	\pagebreak
	\section{Abnahme-Testfälle}

\end{document}